Este trabalho especificou, implementou e avaliou um sistema protótipo \textit{on-premise} para apoio documental ao pré-natal no contexto do SUS, integrando transcrição local, RAG sobre manuais oficiais, LLM executado via Ollama com revisão humana obrigatória e um \textit{gateway} de anonimização de dados.

A pergunta norteadora foi respondida de forma parcial:
a arquitetura mostrou-se tecnicamente viável como base assistiva, auditável e não autônoma, mas os resultados de geração, escopo e latência indicam que o sistema ainda não possui robustez suficiente para uso clínico sem supervisão profissional.
A topologia em contêineres deve ser entendida como decisão de privacidade por desenho, não como comprovação de isolamento absoluto, pois não foram realizados testes específicos de segurança de rede.

Na avaliação experimental, o módulo \textit{RAG} alcançou \textit{Hit@6} global de 78,0\% e MRR médio de 0,4807, mas o \textit{phrase\_recall} médio de 47,1\%, indicando capacidade razoável de recuperar documentos relevantes, embora o \textit{phrase\_recall} mostre que a evidência recuperada nem sempre contém o trecho necessário para satisfazer integralmente o critério de resposta.
O resultado mais expressivo foi observado nos itens classificados como difíceis, para os quais o módulo RAG atingiu Hit@6 de 95,7\%, sugerindo potencial utilidade em cenários menos rotineiros e com maior especificidade protocolar.

O módulo de privacidade não apresentou vazamento no conjunto sintético de 76 entidades avaliadas, embora esse resultado não substitua testes com fala espontânea, abreviações, erros de transcrição e maior variação linguística.
A latência média de 15,59 segundos no fluxo local com \textit{RAG} evidencia que a inferência local e os requisitos de hardware ainda limitam a adoção de baixo custo em produção.
Por fim, a revisão humana obrigatória foi validada como restrição de projeto, pois os resultados reforçam que o sistema deve permanecer assistivo, auditável e não autônomo.
A aprovação de quatro dos cinco casos funcionais do roteiro (GT01-GT05) indica coerência funcional inicial, enquanto a falha no GT05 evidencia a necessidade de barreiras de escopo antes da geração via \textit{LLM}.

Como contribuição, o estudo entrega uma arquitetura replicável, um fluxo de privacidade alinhado aos princípios da LGPD e um benchmark inicial reprodutível para avaliação técnica de recuperação, geração, mascaramento de PII, escopo e latência em documentos de pré-natal.
Os resultados evidenciam os trade-offs entre execução local, custo computacional e segurança operacional, indicando que avanços em modelos, quantização, aceleração em hardware e técnicas de RAG podem tornar esse tipo de sistema mais viável em ambientes reais.
